\subsection{Maturity structure and repeated shocks}\label{app:maturity}
The one-period model concentrates all newly promised service at the next payment date. With longer maturity, the sovereign instead promises a sequence of coupons and principal, and the investor prices survival along that sequence. For a bond with currency coupons $K_{t+j}$, principal $F$ due after $N$ periods, default date $\tau$, and currency recovery $\mathcal R_{t+j}$, a risk-neutral pricing representation is
\begin{align}\label{eq:maturity_price}
 P_t=\E_t\bigg[&\sum_{j=1}^{N}\frac{K_{t+j}\ind\{\tau>t+j\}}{R_f^j}
 +\frac{F\ind\{\tau>t+N\}}{R_f^N}\\
 &+\sum_{j=1}^{N}\frac{\mathcal R_{t+j}\ind\{\tau=t+j\}}{R_f^j}\bigg].\nonumber
\end{align}
This expression does not factor the survival probabilities into independent annual hazards. Future default decisions depend on outstanding claims, refinancing opportunities, and the resources realized by the issuer. The maturity dimension is central to quantitative sovereign-debt models such as \textcite{ChatterjeeEyigungor2012}; introducing it here would therefore require a debt-state problem rather than simply adding more coupons to the current proceeds curve.

For climate applications, maturity matters because immediate disruption and the recovery of resources need not occur on the same schedule. A short event can damage a port or grid whose reconstruction takes several years, while repeated hazards can arrive before the previous service suspension has been repaid. A maturity extension can shift service into a recovery phase, but can also leave obligations concentrated at a later date that remains exposed to another shock. The relevant benefit depends on this joint timing, not just on the number of periods by which principal is moved.

The stationary deferral model in Appendix~\ref{app:deferral} isolates one part of this issue. It adds a second payment date and an inherited obligation but retains i.i.d. growth and a fixed one-period postponement. It does not impose an estimated reconstruction process or allow all outstanding vintages to be restructured repeatedly. A richer maturity model would track the full promised-service schedule and price it under the same default event. Such a model could assess whether smoothing service improves fiscal space when shocks persist, a question that the one-period growth wedge alone cannot answer.

\subsection{State-dependent surplus capacity and adaptation}\label{app:adaptation}
The main GDP benchmark benefits from observing the variable that determines capacity. Once the feasible surplus responds separately to climate or logistics stress, two states with the same output can have different resources available for service. Disaster response can absorb public funds even when measured production has recovered, and a disruption to critical imports can change spending needs independently of its contemporaneous effect on aggregate GDP. In that setting, the relevant state is the joint realization of growth, fiscal capacity, and any recovery or insurance inflows.

This extension changes what an informative trigger means. A GDP clause may continue to capture broad macroeconomic stress, while a verified disaster indicator may capture a fiscal obligation that output alone misses. The ranking depends on the covariance of each index with total repayment resources. It cannot be deduced by inserting a declining surplus into a model in which the GDP ranking relies on a constant surplus. Estimating or calibrating the joint capacity process is therefore a prerequisite for claiming that fiscal fatigue strengthens one design relative to another.

Adaptation adds an intertemporal choice to these resources. An investment can lower future disaster exposure or reduce the loss conditional on an event, but its financing uses present resources or requires additional issuance. The immediate effect on borrowing capacity and the later effect on the growth distribution must both enter a feasible public budget. A verified adaptation stock could be relevant to contractual eligibility or a resilience-linked payment rule, provided that progress and the corresponding reduction in exposure can be measured. The present simulations neither assign a return to adaptation spending nor treat a certification label as an automatic improvement in capacity.

A carbon budget has a different role. A constraint such as $\sum_t E_t\le\bar E$ restricts feasible emissions paths and hence the evolution of chronic damages. By itself, it does not establish a finite proceeds maximum or determine the value of a contingent instrument. Those conclusions require the repayment-resource constraint and a pricing rule. A complete emissions-and-adaptation problem would connect these choices through investment costs, output adjustment, damages, and the liability schedule. Stating these separate requirements avoids interpreting contract design as a substitute for the real-resource problem of mitigation and adaptation.

\subsection{Discount factors, the safe rate, and volatility}\label{app:discounting}
General discount-factor pricing separates expected repayment from the price of its exposure to aggregate states. With zero recovery and $S$ denoting survival, \eqref{eq:sdf} can be written
\begin{equation}\label{eq:sdf_covariance}
 q_H=\E_t[M_{t+1}]\E_t[H S]
       +\operatorname{Cov}_t(M_{t+1},H S).
\end{equation}
The main risk-neutral exercise sets the covariance term to zero and $\E[M]=1/R_f$. If repayments are low precisely in states with a high marginal value of funds to investors, the covariance can instead reduce the price of a claim with a high expected payment. Global climate shocks can be relevant to this channel because a common exposure need not be diversifiable across sovereign issuers. A country-specific default probability alone is insufficient to identify that systematic compensation.

The safe rate itself can also change as the macroeconomic environment changes. A lower discount rate raises the borrowing factor for a given repayment distribution; greater risk or a change in aggregate saving can move the benchmark in the opposite direction. \textcite{Blanchard2019} provides the broader public-debt context for low safe interest rates, without identifying the rate used in these illustrations. Endogenizing it here would require a model of the investors or policy environment that determines the discount factor. The potentially large stationary multiplier makes that omission quantitatively relevant, but it does not supply a sign for the equilibrium response.

State-dependent volatility similarly changes more than a label on the baseline distribution. A change in dispersion can alter the mass near the default threshold, the probability of receiving full payment above a trigger, and the optimizing face. For a lognormal component, changing variance at a fixed log drift also changes mean output. A sensitivity exercise intended to hold mean resources constant would have to adjust the drift accordingly. Neither a higher variance nor a downward stochastic shift is enough to establish the relative value of a step clause without additional information about its trigger and the debt it supports.

\subsection{Private equity and venture capital transmission}\label{app:pevc}
Asset-market transmission can be expressed as an extension of the same state-pricing framework. Let $CF_{t+j}$ denote a private investment's distributable cash flow and $M_{t,t+j}$ a discount factor over $j$ periods. Its value is represented by
\begin{equation}\label{eq:private_value}
 V_t=\E_t\left[\sum_{j=1}^{N}M_{t,t+j}CF_{t+j}\right].
\end{equation}
Climate and logistics shocks can affect both objects. Damage to production, interruption of inputs, and changes in expected demand can lower cash flows. Changes in aggregate risk and financing conditions can change the discount factor applied to those flows. This representation makes the two channels distinct without attributing a measured private-asset effect to the sovereign-debt simulations.

One possible reduced-form discount-rate channel would let a private issuer's required return include a nonnegative exposure loading on a sovereign financing premium. The loading would reflect the importance of domestic financial intermediation, public infrastructure, or the local currency environment for that issuer. It would have to be estimated rather than read from the coefficient on logistics in the sovereign growth process. Evidence that sovereign spreads share global factors, as in \textcite{LongstaffEtAl2011}, is relevant to the common-state interpretation, but does not identify a causal pass-through from a particular contingent clause to a private firm's discount rate.

The timing is especially consequential for private equity and venture capital. A large share of value can depend on cash flows or an exit occurring well after the initial investment. A temporary interruption to operations and a persistent change in the discount applied to future exits can therefore have different valuation implications. The same is true of a change in public financing capacity: it may affect expected infrastructure provision or the availability of domestic finance without changing the physical exposure of the private investment itself. Treating all of these effects as one contemporaneous sovereign-spread coefficient would conceal economically different mechanisms.

A quantitative implementation would require exposure data, cash-flow or exit assumptions, and a stochastic discount factor consistent with the sovereign-pricing block. It would also need to account for common global shocks that can drive sovereign and private values simultaneously. This paper does not supply those objects, so the appendix provides a transmission framework rather than a valuation estimate. Its role is to preserve the connection between climate-sensitive public balance sheets and private capital while keeping the two demonstrated contributions focused on fiscal limits and repayment design.
